\documentclass[12pt]{article}
\usepackage[utf8]{inputenc}
\usepackage[T1]{fontenc}
\usepackage[margin=1in]{geometry}
\usepackage{setspace}
\usepackage{xcolor}
\usepackage[colorlinks=true, urlcolor=blue!60!black, linkcolor=blue!60!black]{hyperref}
\setlength{\parindent}{0em}
\setlength{\parskip}{0.45em}
\renewcommand{\baselinestretch}{1.0}
\begin{document}
\begin{center}
{\Large \textit{Tennessee v. Lane}, 541 U.S. 509 (2004)\par}
\vspace{0.3em}
{\normalsize Case Brief --- One-Page Edition\par}
\end{center}

\textbf{Memory Jogger.} ADA Title~II, sovereign immunity, and access to the courts (\href{https://www.law.cornell.edu/supremecourt/text/541/509}{541 U.S. 509}).

\textbf{Facts.} George Lane, a paraplegic, crawled up two flights of stairs to reach a Tennessee courtroom with no elevator; when he refused to do so again, he was jailed for failure to appear. He and court reporter Beverly Jones sued the State for damages under Title~II.

\textbf{Procedural History.} The district court denied Tennessee's immunity motion; the Sixth Circuit affirmed.

\textbf{Votes.} Affirmed, 5--4. Stevens wrote for the Court; Souter and Ginsburg concurred; Rehnquist, Scalia, and Thomas dissented.

\textbf{Holding.} Assessed as applied to the category of rights at stake, Title~II is congruent and proportional Section~5 legislation, validly abrogating immunity in cases implicating the fundamental right of access to the courts.

\textbf{Reasoning.} Unlike Title~I in \textit{Garrett}, Title~II enforces fundamental rights subject to heightened scrutiny, and the record documented inaccessible courthouses; its reasonable-modification remedy is proportional. Ginsburg described accommodation as equal citizenship. Rehnquist said the majority rigged the test with a courthouse-only statute; Scalia rejected the test as flabby.

\textbf{Examples.} \emph{Same-Side:} a deaf parent denied an interpreter in court. \emph{Opposite-Side:} hockey-rink seating. \emph{Fence-Sitter:} extended exam time at a state university.

\textbf{Critique.} Textualists reject prophylactic enforcement; disability scholars fault the piecemeal holding.

\textbf{Key Quotation.} ``For these reasons, we conclude that Title II, as it applies to the class of cases implicating the fundamental right of access to the courts, constitutes a valid exercise of Congress' \S~5 authority to enforce the guarantees of the Fourteenth Amendment.'' --- Justice Stevens, Opinion of the Court, 541 U.S. at 533--534.

\end{document}
